% Vietnamese translation for OI Public Library
% Source-PDF: IOI中国国家候选队论文/2008/Day2/5.陈瑜希《Pólya计数法的应用》/pálya计数法的应用.pdf
\documentclass[11pt]{article}
\usepackage{oipl-vn}

\title{Ứng Dụng Của Phương Pháp Đếm Pólya}
\author{Chen Yuxi}
\date{2008}

\begin{document}
\maketitle

\origtitle{Applications of Pólya counting}
\sourcepath{2008/Day2/5-Chen-Yuxi/polya-counting-applications.pdf}

\begin{abstract}
Trong các kỳ thi tin học, ta gặp rất nhiều bài toán đếm. Nhiều bài trông có
vẻ khó, nhưng sau khi nắm vững phương pháp đếm Pólya thì có thể giải quyết
khá nhẹ nhàng. Bài viết bắt đầu từ một ví dụ trong thi tin học, trước hết chỉ
ra khó khăn khi dùng phương pháp đếm thông thường, sau đó giới thiệu các khái
niệm cơ bản và tính chất của nhóm, hoán vị, nhóm hoán vị. Trên nền đó, bài
viết đưa vào định lý Burnside, suy ra phương pháp đếm Pólya và chứng minh nó.
Cuối cùng, qua một vài ví dụ, bài viết minh họa ứng dụng của phương pháp đếm
Pólya trong thi tin học và đưa ra tổng kết.
\end{abstract}

\noindent\term{Từ khóa:} định lý Burnside; phương pháp đếm Pólya.

\section{Đặt vấn đề}

\paragraph{Ví dụ 1. He's Circles, SGU 294.}
Cho một vòng tròn gồm \(N\) vị trí, trên đó viết các ký tự `X' và `E'. Hỏi có
bao nhiêu vòng khác nhau về bản chất? Ở đây \(N\le 200000\).

\paragraph{Phân tích.}
Vì đối tượng là một vòng, nhiều phương án vốn khác nhau trước khi quay sẽ trở
thành cùng một phương án sau khi quay. Nếu chỉ dùng nguyên lý nhân để tính thì
không thể loại bỏ các phương án trùng nhau này. Nếu muốn dùng liệt kê, ta phải
liệt kê tất cả phương án. Số lượng liệt kê sẽ không thấp hơn số vòng khác nhau
về bản chất; thực tế số vòng như vậy có cấp \(2^n\). Với quy mô
\(N=200000\), đáp án có hơn sáu vạn chữ số, nên liệt kê rõ ràng không khả thi.

Trong tổ hợp có một phương pháp đếm có thể xử lý rất tốt loại bài toán này.

\section{Kiến thức chuẩn bị}

Dưới đây ta giới thiệu một công cụ đếm quan trọng: phương pháp đếm Pólya. Để
hiểu phương pháp này, trước hết cần xem các khái niệm mà nó sử dụng.

\paragraph{Nhóm.}
Cho tập \(G=\{a,b,c,\ldots\}\) và một phép toán hai ngôi trên \(G\). Nếu thỏa
mãn:
\begin{enumerate}[label=(\alph*)]
  \item \term{Tính đóng}: với mọi \(a,b\in G\), tồn tại \(c\in G\) sao cho
  \(a*b=c\).
  \item \term{Luật kết hợp}: với mọi \(a,b,c\in G\),
  \[
    (a*b)*c=a*(b*c).
  \]
  \item \term{Phần tử đơn vị}: tồn tại \(e\in G\) sao cho với mọi \(a\in G\),
  \[
    a*e=e*a=a.
  \]
  \item \term{Phần tử nghịch đảo}: với mọi \(a\in G\), tồn tại \(b\in G\) sao
  cho
  \[
    a*b=b*a=e,
  \]
  và ký hiệu \(b=a^{-1}\).
\end{enumerate}
thì tập \(G\) dưới phép toán \(*\) được gọi là một nhóm, gọi tắt là \(G\) là
nhóm. Thông thường viết gọn \(a*b\) thành \(ab\).

\paragraph{Hoán vị.}
Một hoán vị giữa \(n\) phần tử \(1,2,\ldots,n\) được viết là
\[
  \begin{pmatrix}
    1 & 2 & \cdots & n\\
    a_1 & a_2 & \cdots & a_n
  \end{pmatrix}.
\]
Nó biểu thị rằng \(1\) được thay bởi một số \(a_1\) trong \(1..n\), \(2\) được
thay bởi một số \(a_2\) trong \(1..n\), cứ thế cho tới \(n\) được thay bởi
\(a_n\); đồng thời \(a_1,a_2,\ldots,a_n\) đôi một khác nhau.

\paragraph{Nhóm hoán vị.}
Các phần tử của nhóm hoán vị là hoán vị, còn phép toán là phép hợp thành hoán
vị. Ví dụ:
\[
  \begin{pmatrix}
    1&2&3&4\\
    3&1&2&4
  \end{pmatrix}
  \begin{pmatrix}
    1&2&3&4\\
    4&3&2&1
  \end{pmatrix}
  =
  \begin{pmatrix}
    1&2&3&4\\
    2&4&3&1
  \end{pmatrix}.
\]
Có thể kiểm tra rằng nhóm hoán vị thỏa mãn bốn điều kiện của nhóm.

Trong ví dụ 1, nhóm hoán vị là
\[
  G=\{\text{quay }0\text{ ô},\text{ quay }1\text{ ô},\ldots,
      \text{ quay }(n-1)\text{ ô}\}.
\]
Các hoán vị tương ứng có dạng
\[
  \begin{pmatrix}
    1&2&\cdots&n\\
    1&2&\cdots&n
  \end{pmatrix},\quad
  \begin{pmatrix}
    1&2&\cdots&n\\
    2&3&\cdots&1
  \end{pmatrix},\quad
  \begin{pmatrix}
    1&2&\cdots&n\\
    3&4&\cdots&2
  \end{pmatrix},\quad\ldots
\]
\[
  \begin{pmatrix}
    1&2&\cdots&n\\
    n&1&\cdots&n-1
  \end{pmatrix}.
\]

\section{Bổ đề Burnside}

\subsection{Giới thiệu}

Sau đây ta giới thiệu bổ đề cần dùng cho phương pháp đếm Pólya: định lý
Burnside. Ký hiệu \(D(a_j)\) là số phần tử bất biến dưới hoán vị \(a_j\), và
\(L\) là số phương án khác nhau về bản chất. Khi đó:
\[
  L=\frac{1}{|G|}\sum_{j=1}^{s}D(a_j).
\]

Trong ví dụ 1, với \(N=4\), có tất cả bốn hoán vị:
\[
  \begin{pmatrix}
    1&2&3&4\\
    1&2&3&4
  \end{pmatrix},\quad
  \begin{pmatrix}
    1&2&3&4\\
    2&3&4&1
  \end{pmatrix},\quad
  \begin{pmatrix}
    1&2&3&4\\
    3&4&1&2
  \end{pmatrix},\quad
  \begin{pmatrix}
    1&2&3&4\\
    4&1&2&3
  \end{pmatrix}.
\]

Mọi phương án đều bất biến dưới hoán vị \(a_1\), nên \(D(a_1)=16\).
Hai phương án \(\texttt{XXXX}\) và \(\texttt{EEEE}\) bất biến dưới \(a_2\),
nên \(D(a_2)=2\). Bốn phương án \(\texttt{XXXX}\), \(\texttt{EEEE}\),
\(\texttt{XEXE}\), \(\texttt{EXEX}\) bất biến dưới \(a_3\), nên
\(D(a_3)=4\). Hai phương án \(\texttt{XXXX}\) và \(\texttt{EEEE}\) bất biến
dưới \(a_4\), nên \(D(a_4)=2\). Do đó
\[
  L=\frac{1}{4}(16+2+4+2)=6.
\]

\subsection{Chứng minh}

Để chứng minh định lý Burnside, ta cần một hệ quả sau. Gọi \(c\) là một cách
tô màu. Số cách tô màu tương đương với \(c\) bằng số hoán vị trong \(G\) chia
cho số hoán vị trong hạt nhân ổn định của \(c\).

\paragraph{Định lý 1.}
Với mỗi cách tô màu \(c\), hạt nhân ổn định \(G(c)\) của \(c\) là một nhóm
hoán vị. Hơn nữa, với mọi hoán vị \(f,g\in G\),
\[
  g*c=f*c
  \quad\Longleftrightarrow\quad
  f^{-1}\circ g\in G(c).
\]

\paragraph{Chứng minh.}
Nếu \(f\) và \(g\) đều giữ \(c\) bất biến, thì thực hiện \(f\) rồi \(g\) cũng
giữ \(c\) bất biến, tức \((g\circ f)(c)=c\). Vì vậy, dưới phép hợp thành,
\(G(c)\) có tính đóng. Rõ ràng phần tử đơn vị giữ mọi cách tô màu bất biến.
Nếu \(f\) giữ \(c\) bất biến thì \(f^{-1}\) cũng giữ \(c\) bất biến, nên
\(G(c)\) đóng với phép lấy nghịch đảo. Do thỏa mãn mọi tính chất trong định
nghĩa nhóm hoán vị, \(G(c)\) là một nhóm hoán vị.

Giả sử \(f*c=g*c\). Tác động \(f^{-1}\) lên hai vế, ta được
\[
  (f^{-1}\circ g)*c=c.
\]
Vì thế \(f^{-1}\circ g\) giữ \(c\) bất biến, tức \(f^{-1}\circ g\in G(c)\).
Ngược lại, nếu \(f^{-1}\circ g\in G(c)\), bằng phép tính tương tự suy ra
\[
  f*c=g*c.
\]

Từ định lý 1, ta có thể bắt đầu từ một cách tô màu \(c\) đã biết và xác định
số cách tô màu khác nhau dưới tác động của \(G\).

\paragraph{Hệ quả 2.}
Gọi \(c\) là một cách tô màu. Số cách tô màu tương đương với \(c\), tức kích
thước quỹ đạo của \(c\), bằng số hoán vị trong \(G\) chia cho số hoán vị trong
hạt nhân ổn định của \(c\):
\[
  |\operatorname{Orb}(c)|=\frac{|G|}{|G(c)|}.
\]

\paragraph{Chứng minh.}
Lấy một hoán vị \(f\in G\). Theo định lý 1, các hoán vị \(g\) thỏa
\[
  g*c=f*c
\]
chính là các hoán vị trong tập
\[
  \{f\circ h\mid h\in G(c)\}.
\]
Theo luật khử, từ \(f\circ h=f\circ h'\) suy ra \(h=h'\). Do đó số hoán vị
trong tập trên bằng số hoán vị trong \(G(c)\). Vì vậy, với mỗi hoán vị \(f\),
có đúng \(|G(c)|\) hoán vị tác động lên \(c\) cho cùng một kết quả như \(f\).
Tổng cộng có \(|G|\) hoán vị, nên số cách tô màu tương đương với \(c\) bằng
\[
  \frac{|G|}{|G(c)|}.
\]

Đã có hệ quả này, việc chứng minh định lý Burnside chỉ còn là áp dụng kỹ
thuật đếm theo hai cách. Ta đếm số cặp \((f,c)\) sao cho \(f\) giữ \(c\) bất
biến, tức \(f*c=c\). Cách đếm thứ nhất là xét từng hoán vị \(f\in G\), đếm số
cách tô màu được \(f\) giữ bất biến, rồi cộng lại. Vì \(D(f)\) là số cách tô
màu bất biến dưới \(f\), cách đếm này cho
\[
  \sum_{f\in G}D(f).
\]

Cách đếm thứ hai là xét từng cách tô màu \(c\), đếm số hoán vị \(f\) thỏa
\(f*c=c\), rồi cộng lại. Với mỗi \(c\), tập các \(f\) như vậy chính là hạt
nhân ổn định \(G(c)\). Do đó cách đếm này cho
\[
  \sum_c |G(c)|.
\]

Nếu phân loại các cách tô màu theo lớp tương đương, tổng trên có thể rút gọn.
Trong cùng một lớp tương đương, mỗi cách tô màu đóng góp cùng một lượng. Với
mỗi lớp tương đương chứa \(c\), theo hệ quả 2, tổng đóng góp của lớp đó là
\[
  |\operatorname{Orb}(c)|\,|G(c)|
  =\frac{|G|}{|G(c)|}|G(c)|
  =|G|.
\]
Số lớp tương đương chính là số cách tô màu không tương đương \(L\), nên
\[
  \sum_c |G(c)|=L|G|.
\]
Suy ra
\[
  L=\frac{1}{|G|}\sum_{f\in G}D(f)
   =\frac{1}{|G|}\sum_{j=1}^{s}D(a_j).
\]

\section{Phương pháp đếm Pólya}

\subsection{Giới thiệu}

Ta nhận thấy việc tính các giá trị \(D(a_j)\) không hề dễ. Nếu dùng tìm kiếm,
độ phức tạp tổng thể là \(O(nsp)\), trong đó \(n\) là số phần tử, \(s\) là số
hoán vị, \(p\) là số ô; ở đây quy mô của \(n\) rất lớn. Bước tiếp theo là tìm
một cách đơn giản để tính \(D(a_j)\). Trước hết giới thiệu khái niệm chu trình.

\paragraph{Chu trình.}
Ký hiệu
\[
  (a_1a_2\cdots a_n)
\]
là một chu trình bậc \(n\). Mỗi hoán vị đều có thể viết thành tích của một số
chu trình rời nhau. Hai chu trình \((a_1a_2\cdots a_n)\) và
\((b_1b_2\cdots b_m)\) rời nhau nghĩa là \(a_i\ne b_j\) với mọi \(i,j\). Ví
dụ:
\[
  \begin{pmatrix}
    1&2&3&4&5\\
    3&5&1&4&2
  \end{pmatrix}
  =(13)(25)(4).
\]
Cách biểu diễn như vậy là duy nhất. Số chu trình của một hoán vị là số chu
trình trong biểu diễn trên. Ví dụ \((13)(25)(4)\) có số chu trình bằng \(3\).

Giả sử \(G\) là một nhóm hoán vị trên \(p\) đối tượng. Dùng \(m\) màu để tô
\(p\) đối tượng đó. Số phương án tô màu khác nhau là
\[
  L=\frac{1}{|G|}
    \left(m^{c(g_1)}+m^{c(g_2)}+\cdots+m^{c(g_s)}\right),
\]
trong đó \(G=\{g_1,\ldots,g_s\}\), còn \(c(g_i)\) là số chu trình của hoán vị
\(g_i\), với \(i=1,\ldots,s\).

Trong ví dụ 1, với vòng có \(N=4\), ta đánh số:
\[
\begin{array}{cc}
1&2\\
4&3
\end{array}
\]
Dựng nhóm hoán vị \(G'=\{g_1,g_2,g_3,g_4\}\), \(|G'|=4\), và gọi \(c(g_i)\)
là số chu trình của \(g_i\). Dưới tác động của \(G'\):
\[
\begin{array}{lll}
g_1\text{ biểu thị quay }0^\circ,   & g_1=(1)(2)(3)(4), & c(g_1)=4,\\
g_2\text{ biểu thị quay }90^\circ,  & g_2=(4\,3\,2\,1), & c(g_2)=1,\\
g_3\text{ biểu thị quay }180^\circ, & g_3=(1\,3)(2\,4), & c(g_3)=2,\\
g_4\text{ biểu thị quay }270^\circ, & g_4=(1\,2\,3\,4), & c(g_4)=1.
\end{array}
\]
Do đó:
\[
  2^{c(g_1)}=2^4=16=D(a_1),\qquad
  2^{c(g_2)}=2^1=2=D(a_2),
\]
\[
  2^{c(g_3)}=2^2=4=D(a_3),\qquad
  2^{c(g_4)}=2^1=2=D(a_4).
\]

Đây chính là định lý Pólya. Ta thấy dùng định lý Pólya có độ phức tạp
\(O(s\times p)\), trong đó \(s\) là số hoán vị và \(p\) là số ô. So với bổ đề
Burnside ở dạng trực tiếp, đây lại là một cải tiến lớn, ưu thế rất rõ ràng.
Định lý Pólya khai thác đầy đủ quan hệ nội tại của đối tượng nghiên cứu, tổng
kết quy luật, bỏ đi nhiều tìm kiếm mù quáng không cần thiết và hạ độ phức tạp
của loại bài toán này xuống mức rất thấp.

\subsection{Chứng minh}

Muốn thu được một phương án ổn định dưới một hoán vị, ta phải tô mọi phần tử
trong cùng một chu trình bằng cùng một màu. Vì vậy
\[
  D(g_i)=m^{c(g_i)}.
\]
Thay biểu thức này vào định lý Burnside, ta thu được công thức Pólya ở trên.

\section{Ứng dụng}

Định lý Pólya có nhiều ví dụ ứng dụng trong thi tin học. Những bài này thường
không thể trực tiếp thay vào công thức để tính, mà cần phân tích kỹ hơn. Sau
đây là một số bài toán đếm trong thi tin học dùng định lý Pólya.

\subsection{Ví dụ 2. Cubes, UVA 10601}

Yêu cầu tô màu \(12\) cạnh của một hình lập phương, trong đó số lượng cạnh
dùng mỗi màu đã được cho. Hãy tính tổng số phương án, coi hai phương án là
một nếu chúng giống nhau sau một phép quay.

\paragraph{Phân tích.}
Bài toán yêu cầu đếm cách tô màu hình lập phương, lại xét bất biến dưới phép
quay, nên rất dễ liên tưởng tới phương pháp đếm Pólya.

Một hình lập phương có tất cả \(24\) phép quay. Dựa trên các phép quay khác
nhau này, ta dựng nhóm hoán vị tương ứng trên các cạnh.

Nếu trực tiếp dùng công thức của phương pháp đếm Pólya, ta không thể bảo đảm
số lần sử dụng của mỗi màu khớp với yêu cầu của đề. Vì vậy cần cải tiến.

Nhìn lại quá trình suy ra công thức Pólya: theo định lý Burnside, số phương
án khác nhau về bản chất bằng tổng số phương án ổn định dưới từng hoán vị,
chia cho tổng số hoán vị. Mặt khác, để một phương án ổn định dưới một hoán
vị, mọi phần tử trong cùng một chu trình của hoán vị phải được tô cùng một
màu, nên \(D(g_i)=m^{c(g_i)}\).

Trong bài này cũng vậy: mỗi chu trình của hoán vị phải được tô cùng một màu.
Vì thế, với mỗi hoán vị, chỉ cần tìm kích thước của từng chu trình, rồi dùng
quy hoạch động nén trạng thái để đếm số cách gán cùng một màu cho mỗi chu
trình sao cho tổng số cạnh của từng màu phù hợp với yêu cầu đề bài.

Điểm khó của bài là có ràng buộc về số lượng từng màu, nên không thể trực
tiếp áp dụng công thức. Ta cần nắm nguồn gốc của phương pháp Pólya, phân tích
thêm một bước, rồi mới thu được cách giải thỏa ràng buộc này.

Việc dùng Pólya trực tiếp còn có một khó khăn khác: vì độ phức tạp của phương
pháp Pólya là \(O(s\times p)\), trong đó \(s\) là số hoán vị và \(p\) là số
ô, nên đôi khi số hoán vị rất lớn và việc thay công thức trực tiếp sẽ quá
thời gian. Muốn tối ưu độ phức tạp, ta phải nghĩ cách giảm số hoán vị cần
duyệt.

\subsection{Ví dụ 3. Transportation is fun, SPOJ 419 và SPOJ 422}

Cho một ma trận \(2^a\times 2^b\). Trong bộ nhớ, ma trận được lưu lần lượt
theo từng hàng: trước hết là hàng thứ nhất, rồi hàng thứ hai, v.v.; trong mỗi
hàng cũng lưu từ trái sang phải. Bây giờ ta muốn thu được ma trận chuyển vị,
vẫn lưu theo cùng cách đó, nhưng chỉ được dùng thao tác hoán đổi. Hỏi cần bao
nhiêu bước hoán đổi?

SPOJ 419 có \(100\) bộ dữ liệu vào; SPOJ 422 có \(400000\) bộ dữ liệu vào.

\paragraph{Phân tích.}
Để tiện mô tả, trước hết xét một ví dụ: giả sử \(a=5\), \(b=3\). Xét phần tử
\((12,1)\), biểu diễn nhị phân là \((01010,001)\). Địa chỉ ban đầu của nó là
\(\texttt{01010 001}\), còn địa chỉ mới là \(\texttt{001 01010}\). Có thể
thấy thao tác chuyển vị này thực chất là dịch vòng địa chỉ của mỗi phần tử
sang phải \(b\) bit.

Nếu hoán vị trên địa chỉ có \(k\) chu trình, thì đáp án là
\[
  2^{a+b}-k.
\]
Số chu trình địa chỉ có thể xem là số địa chỉ khác nhau về bản chất sau khi
áp dụng phép ``dịch vòng \(b\) bit'' trên địa chỉ; đại lượng này có thể tính
bằng nguyên lý Pólya.

Vì SPOJ 422 có số bộ dữ liệu lớn hơn nhiều, ta cần tìm thuật toán tốt hơn.
Với \(a\) và \(b\), trước hết tính \(g=\gcd(a,b)\). Khi đó \(k\) có thể viết
thành
\[
  k=\frac{g}{a+b}
    \sum_{i=1}^{(a+b)/g}
      \left(2^g\right)^{\gcd\left(i,\frac{a+b}{g}\right)}.
\]
Trước hết tiền xử lý các lũy thừa của \(2\), vì Pólya ở đây chỉ dùng một số
lũy thừa của \(2\). Sau đó dùng sàng tìm số nguyên tố, đồng thời tìm thừa số
nguyên tố nhỏ nhất của mỗi hợp số.

Có thể thấy nút thắt là tính các giá trị
\[
  \gcd\left(i,\frac{a+b}{g}\right).
\]
Vì
\[
  \gcd\left(i,\frac{a+b}{g}\right)\mid \frac{a+b}{g},
\]
ta dùng \(f(x)\) biểu thị số lượng \(i\) thỏa
\[
  x=\gcd\left(i,\frac{a+b}{g}\right).
\]

Trước hết xét cách tìm tất cả \(x\) hợp lệ. Vì đã tiền xử lý ``thừa số nguyên
tố nhỏ nhất của mỗi hợp số'', bây giờ có thể tìm mọi thừa số nguyên tố của
\((a+b)/g\) trong thời gian \(O(\log n)\). Có các thừa số nguyên tố rồi thì
việc tìm mọi ước số của nó rất dễ.

Tiếp theo xét cách tính \(f(x)\). Cách làm thật ra rất đơn giản: trong quá
trình suy ra các ước số của \((a+b)/g\), trước hết phân toàn bộ \((a+b)/g\)
giá trị cho \(1\); mỗi lần từ số \(x\) sinh ra \(xp\), chuyển \(1/p\) lượng
của \(f(x)\) sang \(f(xp)\). Như vậy, trong lúc tìm mọi ước số, ta đồng thời
tìm được giá trị của \(f(x)\). Khi đó thay vào công thức phía trên để tính là
rất đơn giản.

Bài này không khó ở bước chuyển thành mô hình đếm Pólya; điểm khó nằm ở việc
tính biểu thức sau khi chuyển đổi. Số lượng hoán vị quá lớn khiến ta không
thể tính số chu trình cho từng hoán vị riêng lẻ, đó là trở ngại chính trong
đếm Pólya. Ví dụ tiếp theo cũng gặp tình huống tương tự, nhưng sau khi dựng
được mô hình Pólya, cách tối ưu thời gian của nó rất khéo léo.

\subsection{Ví dụ 4. Isomorphism, SGU 282}

Một đồ thị tô màu là đồ thị vô hướng đầy đủ, và mỗi cạnh có thể được tô bằng
một trong \(M\) màu. Hai đồ thị tô màu là đẳng cấu khi và chỉ khi có thể đổi
nhãn các đỉnh của một đồ thị để hai đồ thị hoàn toàn giống nhau. Hỏi với
\(N\) đỉnh và \(M\) màu, có bao nhiêu đồ thị tô màu khác nhau về bản chất,
tức đôi một không đẳng cấu, lấy modulo số nguyên tố \(P\).

\[
  1\le N\le 53,\qquad 1\le M\le 1000,\qquad N<P\le 10^9,
\]
giới hạn thời gian là \(10\) giây. Bài này là đề tuyển chọn đội tuyển tỉnh
Thượng Hải và Giang Tô năm 2006.

\paragraph{Phân tích.}
Bài toán dạng này phù hợp với phạm vi áp dụng của Pólya.

Trong bài này, các đối tượng của nhóm hoán vị là \(\binom{N}{2}\) cạnh, số
màu là \(M\), còn \(G\) là nhóm các hoán vị trên cạnh được sinh ra bởi hoán
vị các đỉnh.

Đồng thời, ta biết nếu dùng định lý Pólya trực tiếp thì có thể giải trong
thời gian \(O(ps)\), trong đó \(s=|G|\). Nhưng với bài này, \(|G|=N!\). Khi
\(N=53\), con số này cực kỳ lớn, nên cần phân tích sâu hơn.

Xét một hoán vị trên đỉnh \(f:i\mapsto P[i]\). Hoán vị tương ứng trên cạnh là
\[
  f':(i,j)\mapsto (P[i],P[j]).
\]
Ta xét quan hệ giữa số chu trình \(c(f)\) của hoán vị trên đỉnh và số chu
trình \(c(f')\) của hoán vị trên cạnh. Có thể thu được các kết luận sau.

Giả sử đỉnh \(i\) và đỉnh \(j\) cùng thuộc một chu trình dài \(L\). Khi đó,
các cạnh \((i,j)\) tương ứng tạo thành số chu trình bằng
\[
  \left\lfloor\frac{L}{2}\right\rfloor.
\]

Giả sử đỉnh \(i\) và đỉnh \(j\) lần lượt thuộc hai chu trình khác nhau có độ
dài \(L_1\) và \(L_2\). Khi đó, các cạnh \((i,j)\) tương ứng tạo thành số chu
trình bằng \(\gcd(L_1,L_2)\).

Không mất tính tổng quát, giả sử
\[
  L_1\ge L_2\ge\cdots\ge L_m>0,
  \qquad L_1+L_2+\cdots+L_m=N.
\]
Như vậy \(L_1,L_2,\ldots,L_m\) tạo thành một phân hoạch của \(N\). Theo phân
tích trên, mọi hoán vị ứng với cùng một phân hoạch có cùng số chu trình trên
cạnh; đó là một hàm theo \(L_1,L_2,\ldots,L_m\):
\[
  T_1=\sum_{i=1}^{m}\left\lfloor\frac{L_i}{2}\right\rfloor
      +\sum_{1\le i<j\le m}\gcd(L_i,L_j).
\]

Vì \(N\le 53\), tổng số phân hoạch của \(N\) không lớn. Ta có thể dùng quay
lui để liệt kê mọi phân hoạch của \(N\), rồi với mỗi phân hoạch, tính số hoán
vị tương ứng với phân hoạch đó và số chu trình của mỗi hoán vị.

Số chu trình đã được thảo luận ở trên. Còn số hoán vị là bao nhiêu?

Giả sử đã xác định
\[
  L_1\ge L_2\ge\cdots\ge L_m>0.
\]
Tiếp theo là đưa \(N\) đỉnh \(1,\ldots,N\) vào \(m\) chu trình sao cho chu
trình thứ \(i\) chứa đúng \(L_i\) đỉnh. Phần này có thể xem là một bài toán
tổ hợp đơn giản, có
\[
  \frac{N!}{L_1!L_2!\cdots L_m!}
\]
cách.

Với mỗi chu trình, sau khi đã xác định đỉnh đầu tiên, các cách sắp xếp còn
lại của các đỉnh là khác nhau và thực ra ứng với các hoán vị khác nhau. Vì
vậy còn phải nhân với
\[
  (L_1-1)!(L_2-1)!\cdots(L_m-1)!.
\]
Nếu có \(L_i=L_{i+1}=\cdots=L_j\), thì mỗi \((j-i+1)!\) phương án lại bị lặp,
nên cần chia cho \((j-i+1)!\).

Do đó tổng số hoán vị là
\[
  \frac{N!}{L_1L_2\cdots L_m\, k_1!k_2!\cdots k_t!},
\]
trong đó \(t\) là số giá trị \(L\) khác nhau, và mỗi giá trị xuất hiện \(k_i\)
lần. Những giá trị này đều có thể tính trong lúc liệt kê phân hoạch.

Vì vậy, đáp án của toàn bộ bài toán có thể được tính từ số chu trình và số
hoán vị của từng loại hoán vị.

Sau đây bàn cách tính cụ thể. Mỗi phân hoạch đóng góp một hạng có dạng
\[
  M^{T_1}\cdot T_2^{-1}\pmod P,
  \qquad
  T_2=L_1L_2\cdots L_m\, k_1!k_2!\cdots k_t!.
\]
Phần thứ nhất là \(M^{T_1}\). Ở đây \(T_1\) không lớn, đủ để lưu bằng
\texttt{longint}, và \(M^{T_1}\bmod P\) có thể tính bằng lũy thừa nhanh trong
thời gian \(O(\log T_1)\); đây là kỹ thuật quen thuộc nên không nhắc lại.

Phần còn lại là \(T_2^{-1}\). \(T_2\) rất lớn, lại là lũy thừa \(-1\); chẳng
lẽ phải phân tích thừa số nguyên tố? Không cần. Hãy chú ý một điều kiện rất
nhỏ nhưng rất quan trọng của đề: \(P\) là số nguyên tố và \(N<P\). Điều kiện
này có tác dụng gì?

Rõ ràng
\[
  L_1,\ldots,L_m,k_1,\ldots,k_t\le N<P,
\]
nên mọi \(L_1,\ldots,L_m,k_1,\ldots,k_t\) đều nguyên tố cùng nhau với \(P\).
Vì vậy \(T_2\) cũng nguyên tố cùng nhau với \(P\). Mặt khác, vì \(P\) là số
nguyên tố, theo kiến thức số học:
\[
  T_2^{P-1}\equiv 1\pmod P,
\]
\[
  T_2^{-1}\equiv T_2^{P-2}\pmod P.
\]
Do đó có thể chuyển việc tính \(T_2^{-1}\) thành tính \(T_2^{P-2}\), cũng
giống như cách tính \(M^{T_1}\).

Bài này gặp cùng khó khăn với bài trước: số lượng hoán vị quá lớn khiến không
thể tính số chu trình cho từng hoán vị. Cách giải là tìm ra các hoán vị tương
tự nhau trong nhóm hoán vị và không tính lặp lại. Phương pháp loại bỏ phép
tính dư thừa này thường xuất hiện trong các bài đếm Pólya. Sau khi dùng tư
tưởng đó, việc tính số lượng hoán vị của mỗi lớp tương tự cũng đòi hỏi nền
tảng toán học vững chắc.

\section{Tổng kết}

Bài viết bắt đầu từ một ví dụ trong thi tin học, giới thiệu phương pháp đếm
Pólya, tư tưởng toán học ẩn trong nó và ứng dụng của nó trong thi tin học.
Phương pháp đếm Pólya không chỉ giải được nhiều bài toán đếm; quá trình chứng
minh của nó cũng rất thú vị. Muốn sử dụng linh hoạt phương pháp đếm Pólya, ta
không chỉ cần nắm vững tính chất của loại bài toán này, mà còn cần nền tảng
toán học chắc và năng lực phân tích vấn đề. Một lần nữa: phương pháp toán học
là công cụ để giải bài toán, còn năng lực phân tích vấn đề là nguồn gốc của
thuật toán.

\section*{Tài liệu tham khảo}

\begin{enumerate}
  \item Richard A. Brualdi, \textit{Introductory Combinatorics}. Bản dịch
  tiếng Trung của Feng Shunxi, Luo Ping, Pei Weidong; Lu Kaicheng và Feng
  Shunxi hiệu đính.
  \item Fu Wenjie, bài luận đội tuyển quốc gia năm 2001.
  \item Bài tập đội tuyển quốc gia các năm 2005--2007.
\end{enumerate}

\end{document}
